% Information Processing in the c-Field: Longitudinal Modes, Density-Dependent
% Logic, and Post-Silicon Computing via Density Field Dynamics
%
% Supporting paper for U.S. Provisional Patent Application No. 63/865,293
% Gary Alcock, 2025--2026
%
\documentclass[aps,prd,twocolumn,superscriptaddress,nofootinbib,floatfix]{revtex4-2}

\usepackage{amsmath,amssymb,amsfonts}
\usepackage{graphicx}
\usepackage{hyperref}
\usepackage{xcolor}
\usepackage{bm}
\usepackage{mathrsfs}
\usepackage{braket}
\usepackage{booktabs}
\usepackage{multirow}
\usepackage{array}
\usepackage{tikz}
\usetikzlibrary{arrows.meta,positioning,decorations.pathmorphing,calc,shapes.geometric}
\usepackage{float}

% Custom commands
\newcommand{\psifield}{\psi}
\newcommand{\cfield}{c_{\text{1w}}}
\newcommand{\cvac}{c}
\newcommand{\nref}{n}
\newcommand{\avec}{\mathbf{a}}
\newcommand{\xvec}{\mathbf{x}}
\newcommand{\kvec}{\mathbf{k}}
\newcommand{\Evec}{\mathbf{E}}
\newcommand{\Bvec}{\mathbf{B}}
\newcommand{\Jvec}{\mathbf{J}}
\newcommand{\Svec}{\mathbf{S}}
\newcommand{\Order}{\mathcal{O}}
\newcommand{\Lag}{\mathcal{L}}
\newcommand{\Ham}{\mathcal{H}}
\newcommand{\calB}{\mathcal{B}}
\newcommand{\calC}{\mathcal{C}}
\newcommand{\calF}{\mathcal{F}}
\newcommand{\calG}{\mathcal{G}}
\newcommand{\calI}{\mathcal{I}}
\newcommand{\calN}{\mathcal{N}}
\newcommand{\dfd}{\text{DFD}}
\newcommand{\snr}{\text{SNR}}
\newcommand{\ee}[1]{\times 10^{#1}}

\begin{document}

\title{Information Processing in the $c$-Field:\\
Longitudinal Modes, Density-Dependent Logic,\\
and Post-Silicon Computing via Density Field Dynamics}

\author{Gary Alcock}
\email{gary@gtacompanies.com}
\affiliation{Independent Researcher, Los Angeles, CA, USA}

\date{\today}

\begin{abstract}
Density Field Dynamics (DFD) posits a scalar refractive-index field
$\psifield(\xvec,t)$ on flat spacetime with $n = e^{\psifield}$,
producing a dynamical one-way speed of light
$\cfield = \cvac\, e^{-\psifield}$.
We demonstrate that spatial and temporal variations of $\psifield$
support \emph{longitudinal} electromagnetic propagation modes absent
in standard vacuum electrodynamics, deriving the full dispersion
relation and showing that these modes propagate as compressional
disturbances of the $c$-field.
We compute the information-theoretic channel capacity of a
$c$-field--modulated communication link, proving that the
density-dependent bandwidth exceeds the conventional Shannon limit
for transverse EM channels under equivalent power constraints when
the $\psifield$-gradient exceeds a critical threshold.
We design a complete computing architecture---the \emph{$c$-field
logic unit} (CLU)---in which Boolean and higher-dimensional logic
operations are performed through interference and phase-shifting of
longitudinal $c$-field modes in density-structured waveguides.
We present prototype specifications for a laboratory $c$-field
modulator/detector system capable of demonstrating the
communication principle, with sensitivity targets informed by the
EM--$\psifield$ coupling parameter $\lambda$ constrained to
$|\lambda - 1| \lesssim 3\times10^{-5}$ by existing cavity
stability data.
Applications to ultra-secure communication, high-density post-silicon
computing, hybrid AI architectures, and post-quantum cryptography
are developed in detail.
\end{abstract}

\maketitle
\tableofcontents

%======================================================================
\section{Introduction}
\label{sec:intro}
%======================================================================

\subsection{Motivation: The information-theoretic frontier}

Modern information technology rests on two pillars: electromagnetic
wave propagation for communication and semiconductor switching for
computation. Both face fundamental limits. Transverse electromagnetic
(TEM) waves require line-of-sight or guided paths, suffer frequency-dependent
attenuation in matter, and are blocked by Faraday shielding.
Silicon computing approaches thermodynamic and quantum-tunneling
barriers as feature sizes shrink below 3~nm~\cite{IRDS2024}.
Quantum computing, while promising exponential speedup for certain
problems, faces decoherence, error-correction overhead, and extreme
cryogenic requirements that limit near-term scalability~\cite{Preskill2018}.

Photonic computing has emerged as an alternative, exploiting the
speed-of-light propagation and natural parallelism of optical
signals~\cite{Shen2017,Feldmann2021}. However, photonic architectures
remain constrained by the same transverse-mode physics of standard
electrodynamics: diffraction limits set minimum feature sizes, and
nonlinear optical elements required for logic operations are
inherently lossy and difficult to integrate~\cite{Miller2017}.

This paper demonstrates that Density Field Dynamics (DFD) opens an
entirely new parameter space for information processing. The DFD
framework, developed in Refs.~\cite{Alcock2025DFD,Alcock2025Strong,
Alcock2025EMCoupling}, establishes that the one-way speed of light
is a dynamical, position-dependent quantity $\cfield = \cvac\,
e^{-\psifield}$ determined by a scalar refractive-index field
$\psifield(\xvec,t)$.
The key consequences for information science are:

\begin{enumerate}
\item \textbf{Longitudinal EM modes.} Gradients $\nabla\psifield
\neq 0$ break the transverse-only constraint of vacuum
electrodynamics, enabling compressional electromagnetic waves
that propagate as density variations of the $c$-field.

\item \textbf{Density-dependent signal pathways.} Because the
refractive index $n = e^{\psifield}$ varies with local energy
density, information can be routed through regions that are opaque
to conventional transverse EM, bypassing attenuation, shielding,
and line-of-sight constraints.

\item \textbf{$c$-field modulation as information carrier.} The
local value of $\cfield$ provides a continuously variable
physical parameter into which information can be encoded at
densities exceeding conventional electromagnetic encoding.
\end{enumerate}

\subsection{Relation to prior work}

The concept of longitudinal electromagnetic waves has a long and
contested history. Tesla's early experiments suggested non-transverse
propagation modes~\cite{Tesla1904}. More recently, theoretical work
on ``scalar waves'' in extended electrodynamics has been pursued by
several groups~\cite{Meyl2003,Monstein2002}, though these claims
remain outside the mainstream consensus.
Within standard physics, longitudinal modes arise naturally in
plasma physics (Langmuir waves)~\cite{Chen1984}, in waveguide
TM modes with nonzero $E_z$ components, and in the near-field
of antennas~\cite{Jackson1998}. The Aharonov-Bohm effect
demonstrates that electromagnetic potentials---including the
scalar potential, which has longitudinal character---carry
physical information beyond the transverse fields~\cite{Aharonov1959}.

DFD provides a \emph{rigorous theoretical framework} within which
longitudinal EM propagation emerges as a natural consequence of
the variable-$c$ structure, with no ad hoc extensions to Maxwell's
equations. The longitudinal mode is not a separate field; it is a
compressional disturbance of the same $\psifield$ field that
governs gravity, lensing, and clock rates.

For computing, the closest precedent is metamaterial-based wave
computing~\cite{Silva2014,Estakhri2019}, where structured dielectric
environments perform mathematical operations on propagating waves.
DFD extends this concept: rather than static metamaterial structures,
the computing medium is the dynamical $\psifield$ field itself,
whose density-dependent response provides intrinsic nonlinearity
for logic operations.

\subsection{Plan of this paper}

Section~\ref{sec:framework} reviews the DFD framework relevant to
electrodynamics and establishes notation.
Section~\ref{sec:longitudinal} derives the longitudinal EM mode
from first principles, obtains its dispersion relation, and
characterizes propagation, attenuation, and coupling to the
transverse sector.
Section~\ref{sec:channel} develops the information-theoretic
channel capacity of $c$-field communication.
Section~\ref{sec:computing} designs the $c$-field logic unit and
develops the full computing architecture.
Section~\ref{sec:applications} presents specific applications.
Section~\ref{sec:prototype} gives the engineering specification
for a prototype demonstration.
Section~\ref{sec:discussion} discusses implications and limitations.

%======================================================================
\section{DFD Electrodynamic Framework}
\label{sec:framework}
%======================================================================

\subsection{Core postulates}

Density Field Dynamics~\cite{Alcock2025DFD} is built on two postulates:

\textbf{P1 (Light).} In a broadband nondispersive window,
electromagnetic waves propagate according to the eikonal of an
effective optical metric
\begin{equation}
d\tilde{s}^2 = -\frac{\cvac^2\,dt^2}{n^2(\xvec,t)} + d\xvec^2,
\qquad n(\xvec,t) = e^{\psifield(\xvec,t)}.
\label{eq:optical_metric}
\end{equation}

\textbf{P2 (Matter).} Test bodies move under the conservative
potential
\begin{equation}
\Phi \equiv -\frac{\cvac^2}{2}\,\psifield, \qquad
\avec = \frac{\cvac^2}{2}\,\nabla\psifield = -\nabla\Phi.
\label{eq:matter_potential}
\end{equation}

The one-way phase velocity of light is
\begin{equation}
\cfield(\xvec,t) = \frac{\cvac}{n} = \cvac\,e^{-\psifield(\xvec,t)},
\label{eq:c_field}
\end{equation}
which is the central dynamical object of this paper. In a
gravitational potential well ($\psifield > 0$), light slows:
$\cfield < \cvac$.

\subsection{The scalar field equation}

The scalar field $\psifield$ satisfies the nonlinear elliptic
equation~\cite{Alcock2025DFD}:
\begin{equation}
\nabla\cdot\!\left[\mu\!\left(\frac{|\nabla\psifield|}{a_\star}\right)
\nabla\psifield\right] = -\frac{8\pi G}{\cvac^2}(\rho - \bar\rho),
\label{eq:field_equation}
\end{equation}
where $\mu(x) = x/(1+x)$ is the response function (derived from
$S^3$ Chern-Simons topology~\cite{Alcock2025DFD}), $a_\star$ is
the characteristic gradient scale related to the MOND acceleration
$a_0 = 2\sqrt{\alpha}\,\cvac H_0$ via $a_\star = 2a_0/\cvac^2$,
$\rho$ is local mass density, and $\bar\rho$ is the cosmic mean.

For the information-processing applications in this paper, we work
in the \emph{laboratory regime} where
$|\nabla\psifield|/a_\star \gg 1$, so $\mu \to 1$ and the field
equation reduces to the Poisson equation:
\begin{equation}
\nabla^2\psifield = -\frac{8\pi G}{\cvac^2}\,\rho.
\label{eq:poisson_limit}
\end{equation}

\subsection{EM--$\psifield$ coupling: The $\lambda$ and $\kappa$ parameters}
\label{subsec:lambda_kappa}

The EM--$\psifield$ interaction is governed by two dimensionless
parameters~\cite{Alcock2025EMCoupling,Alcock2025DFD}:

\textbf{The $\lambda$ parameter.} Controls whether EM fields
can source $\psifield$ oscillations. When $\lambda = 1$, EM
probes the optical metric $n = e^{\psifield}$ but does not pump
$\psifield$. When $|\lambda - 1| \neq 0$, EM fields actively
drive $\psifield$ oscillations. The mode equation for a single
laboratory $\psifield$ mode $q(t)$ with natural frequency
$\Omega_\psi$ and damping $\gamma_\psi$ is~\cite{Alcock2025EMCoupling}:
\begin{equation}
\ddot{q} + 2\gamma_\psi\dot{q} + \Omega_\psi^2 q
= \frac{(\lambda-1)}{M_\psi}\int u(\mathbf{r})\,
\Xi(\mathbf{r},t)\,d^3r + \alpha\,U(t)\,q,
\label{eq:mode_equation}
\end{equation}
where $\Xi(\mathbf{r},t) \equiv -\tfrac{1}{2}\,e^{-2\psifield_0}
(B^2 - E^2/\cvac^2)$ is the EM stress tensor trace, $u(\mathbf{r})$
is the normalized spatial profile, $M_\psi$ is the effective mode
mass, and $U(t) = U_0[1 + m\cos(2\omega t)]$ is the stored EM
energy with modulation depth $m$.

Existing cavity stability constrains~\cite{Alcock2025EMCoupling}:
\begin{equation}
|\lambda - 1| \lesssim 3\times 10^{-5}
\quad\text{(accidental bound)}.
\label{eq:lambda_bound}
\end{equation}

\textbf{The $\kappa$ parameter.} Controls the differential
response of electric and magnetic sectors to $\psifield$:
\begin{equation}
\epsilon(\psifield) = \epsilon_0\,n\,e^{+\kappa\psifield},
\qquad
\mu_{\text{vac}}(\psifield) = \mu_0\,n\,e^{-\kappa\psifield},
\label{eq:kappa_split}
\end{equation}
preserving the phase velocity $v_{\text{ph}} = 1/\sqrt{\epsilon\mu}
= \cvac/n = \cfield$. The unified bracket
$\calB \equiv B^2/\mu - \epsilon E^2$ governs energy exchange,
body force $\mathbf{f}_\psi = -(\kappa/2)\,\calB\,\nabla\psifield$,
and $\psifield$ sourcing~\cite{Alcock2025DFD}.

\subsection{Modified Maxwell equations in the $\psifield$-medium}
\label{subsec:maxwell_mod}

In the DFD optical medium with refractive index
$n(\xvec,t) = e^{\psifield(\xvec,t)}$, Maxwell's equations become:
\begin{align}
\nabla\cdot(\epsilon\,\Evec) &= \rho_f,
\label{eq:gauss_E}\\
\nabla\cdot\Bvec &= 0,
\label{eq:gauss_B}\\
\nabla\times\Evec &= -\partial_t\Bvec,
\label{eq:faraday}\\
\nabla\times\!\left(\frac{\Bvec}{\mu_{\text{vac}}}\right)
&= \Jvec_f + \partial_t(\epsilon\,\Evec),
\label{eq:ampere}
\end{align}
where $\epsilon = \epsilon(\psifield)$ and
$\mu_{\text{vac}} = \mu_{\text{vac}}(\psifield)$ are the
position-dependent permittivity and permeability from
Eq.~(\ref{eq:kappa_split}).

The critical observation is that when $\nabla\psifield \neq 0$,
the spatial derivatives of $\epsilon$ and $\mu_{\text{vac}}$
generate source terms that couple the longitudinal and transverse
components of the electromagnetic field. This is the origin of the
longitudinal mode.

%======================================================================
\section{Longitudinal Electromagnetic Modes in DFD}
\label{sec:longitudinal}
%======================================================================

\subsection{Derivation from first principles}

We work with the wave equation derived from the modified Maxwell
equations~(\ref{eq:gauss_E})--(\ref{eq:ampere}).
Taking the curl of Eq.~(\ref{eq:faraday}) and substituting
Eq.~(\ref{eq:ampere}) in a source-free region ($\rho_f = 0$,
$\Jvec_f = 0$):
\begin{equation}
\nabla\times(\nabla\times\Evec) =
-\partial_t\!\left[\nabla\times\Bvec\right]
= -\mu_{\text{vac}}\,\epsilon\,\partial_t^2\Evec
- \mu_{\text{vac}}\,(\partial_t\epsilon)\,\partial_t\Evec.
\label{eq:wave_step1}
\end{equation}

Using the vector identity
$\nabla\times(\nabla\times\Evec) = \nabla(\nabla\cdot\Evec) -
\nabla^2\Evec$ and noting that $\nabla\cdot\Evec \neq 0$ when
$\nabla\epsilon \neq 0$ (from Eq.~(\ref{eq:gauss_E}) with
$\rho_f = 0$):
\begin{equation}
\nabla\cdot(\epsilon\,\Evec) = 0 \;\Rightarrow\;
\nabla\cdot\Evec = -\frac{\nabla\epsilon}{\epsilon}\cdot\Evec
= -(1+\kappa)\,\nabla\psifield\cdot\Evec,
\label{eq:div_E_nonzero}
\end{equation}
where we used
$\nabla\ln\epsilon = (1+\kappa)\nabla\psifield$.

Define the \emph{longitudinal polarization} as the component of
$\Evec$ along $\nabla\psifield$:
\begin{equation}
E_L \equiv \hat{\mathbf{g}}\cdot\Evec, \qquad
\hat{\mathbf{g}} \equiv \frac{\nabla\psifield}{|\nabla\psifield|}.
\label{eq:E_longitudinal}
\end{equation}

The transverse components satisfy
$\Evec_T = \Evec - E_L\,\hat{\mathbf{g}}$.

\subsection{The longitudinal wave equation}

For a background with slowly varying $\psifield$ (WKB regime:
$|\nabla^2\psifield|/|\nabla\psifield|^2 \ll k$ where $k$ is
the EM wavenumber), projecting the wave equation onto
$\hat{\mathbf{g}}$ yields:
\begin{equation}
\boxed{
\frac{\partial^2 E_L}{\partial t^2}
- \cfield^2\,\frac{\partial^2 E_L}{\partial s^2}
+ \Gamma_L\,\frac{\partial E_L}{\partial t}
+ \omega_{\text{cut}}^2\,E_L = \calS_{LT},
}
\label{eq:longitudinal_wave}
\end{equation}
where:
\begin{itemize}
\item $s$ is the coordinate along $\hat{\mathbf{g}}$
(the $\psifield$-gradient direction),

\item $\cfield = \cvac\,e^{-\psifield}$ is the local propagation
speed,

\item $\Gamma_L = (1+\kappa)\,\dot\psifield$ is a damping/growth
rate set by the temporal variation of $\psifield$,

\item $\omega_{\text{cut}}^2 = (1+\kappa)\,\cfield^2\,
\nabla^2\psifield + \Order(|\nabla\psifield|^2)$ is a cutoff
frequency analogous to the plasma frequency in Langmuir waves,

\item $\calS_{LT}$ is a source term representing coupling from
transverse modes, proportional to $(1+\kappa)\,|\nabla\psifield|$
times the transverse field amplitude.
\end{itemize}

\subsection{Dispersion relation}
\label{subsec:dispersion}

For a monochromatic mode $E_L \propto e^{i(ks - \omega t)}$
propagating through a region of approximately uniform
$\psifield$-gradient, the dispersion relation is:
\begin{equation}
\omega^2 = \cfield^2\,k^2 + \omega_{\text{cut}}^2
- i\,\omega\,\Gamma_L.
\label{eq:dispersion}
\end{equation}

In the undamped, above-cutoff regime
($\omega > \omega_{\text{cut}}$, $\Gamma_L \to 0$):
\begin{equation}
k = \frac{1}{\cfield}\sqrt{\omega^2 - \omega_{\text{cut}}^2},
\label{eq:wavenumber}
\end{equation}
giving a group velocity:
\begin{equation}
v_g = \frac{\partial\omega}{\partial k}
= \cfield\sqrt{1 - \frac{\omega_{\text{cut}}^2}{\omega^2}}.
\label{eq:group_velocity}
\end{equation}

The phase velocity exceeds $\cfield$ for
$\omega > \omega_{\text{cut}}$:
\begin{equation}
v_p = \frac{\omega}{k}
= \frac{\cfield}{\sqrt{1 - \omega_{\text{cut}}^2/\omega^2}}.
\label{eq:phase_velocity}
\end{equation}

This is the same dispersive structure as electromagnetic waves in
a plasma, with $\omega_{\text{cut}}$ playing the role of the plasma
frequency. The physical origin is different: here it arises from the
curvature of the $\psifield$ field (i.e., $\nabla^2\psifield$),
not from free charge carriers.

\subsection{Cutoff frequency estimate}

For a laboratory-generated $\psifield$-gradient using a mass
distribution of density $\rho$ and characteristic scale $L$:
\begin{equation}
\omega_{\text{cut}}^2 \sim (1+\kappa)\,\frac{\cvac^2}{n^2}\,
\frac{8\pi G\rho}{\cvac^2}
= (1+\kappa)\,\frac{8\pi G\rho}{n^2}.
\label{eq:cutoff_estimate}
\end{equation}

For terrestrial conditions ($\rho \sim 10^4$~kg/m$^3$,
$n \approx 1$):
\begin{equation}
\omega_{\text{cut}} \sim \sqrt{8\pi \times 6.67\ee{-11}
\times 10^4} \sim 10^{-3}\text{ rad/s},
\label{eq:cutoff_numerical}
\end{equation}
corresponding to a cutoff frequency
$f_{\text{cut}} \sim 10^{-4}$~Hz.

This extremely low cutoff means that \emph{all practical EM
frequencies are far above cutoff}, and the longitudinal mode
propagates with negligible dispersion:
$v_g \approx v_p \approx \cfield$ for $\omega \gg \omega_{\text{cut}}$.

\subsection{Comparison: Longitudinal vs.\ transverse modes}
\label{subsec:comparison}

\begin{table}[h]
\caption{Comparison of transverse and longitudinal EM modes in DFD.}
\label{tab:mode_comparison}
\begin{ruledtabular}
\begin{tabular}{lcc}
Property & Transverse & Longitudinal \\
\hline
Polarization & $\Evec \perp \kvec$ & $\Evec \parallel \nabla\psifield$ \\
Propagation speed & $\cfield$ & $\cfield$ (above cutoff) \\
Dispersion & None & $\omega_{\text{cut}} \sim 10^{-4}$~Hz \\
Requires $\nabla\psi$ & No & Yes \\
Faraday shielding & Blocked & Penetrates \\
Attenuation in matter & Standard & Density-coupled \\
Information carrier & $\Evec$, $\Bvec$ amplitudes & $\delta\psifield$ amplitude \\
\end{tabular}
\end{ruledtabular}
\end{table}

The critical distinction is in the last two rows. Longitudinal modes
propagate as compressions of the $c$-field itself. A Faraday cage
blocks transverse EM by shorting out the electric field at the
conducting boundary. But a $\psifield$-gradient disturbance---being
a variation of the refractive index of spacetime---is not screened
by conducting surfaces. The $\psifield$ field couples to mass-energy,
not to free charges.

\subsection{Longitudinal--transverse mode coupling}

The source term $\calS_{LT}$ in Eq.~(\ref{eq:longitudinal_wave})
provides a mechanism for converting between transverse and
longitudinal modes:
\begin{equation}
\calS_{LT} = (1+\kappa)\,\cfield^2\left[
(\hat{\mathbf{g}}\cdot\nabla)(\nabla\psifield\cdot\Evec_T)
+ |\nabla\psifield|\,\nabla\cdot\Evec_T^{(\psi)}
\right],
\label{eq:mode_coupling}
\end{equation}
where $\Evec_T^{(\psi)}$ denotes the transverse field components
projected onto the $\psifield$-gradient plane.

The coupling efficiency scales as:
\begin{equation}
\eta_{T\to L} \sim (1+\kappa)^2\,
\left(\frac{|\nabla\psifield|\,\lambda_{\text{EM}}}{2\pi}\right)^2,
\label{eq:coupling_efficiency}
\end{equation}
where $\lambda_{\text{EM}}$ is the electromagnetic wavelength.
For laboratory $\psifield$-gradients
($|\nabla\psifield| \sim GM/(c^2 R^2) \sim 10^{-26}$~m$^{-1}$
for Earth surface), this natural coupling is extremely weak:
$\eta_{T\to L} \sim 10^{-50}$. However, \emph{engineered}
$\psifield$-gradients using the EM-$\psifield$ back-reaction
(the $\lambda$ parameter) can enhance this dramatically, as we
show in Sec.~\ref{sec:prototype}.

\subsection{Energy and momentum of longitudinal modes}

The energy density of the longitudinal mode is:
\begin{equation}
u_L = \frac{1}{2}\left[\epsilon\,E_L^2
+ \frac{1}{\mu_{\text{vac}}}\,B_L^2\right]
+ \frac{1}{2}\,(1+\kappa)\,\epsilon\,E_L^2\,
\frac{\nabla^2\psifield}{k^2},
\label{eq:energy_density_L}
\end{equation}
where the second term is the ``medium contribution'' from the
$\psifield$ gradient. The associated Poynting-like flux is:
\begin{equation}
\Svec_L = \frac{E_L^2}{2}\,\sqrt{\frac{\epsilon}{\mu_{\text{vac}}}}
\,v_g\,\hat{\mathbf{g}},
\label{eq:poynting_L}
\end{equation}
directed along the $\psifield$-gradient, confirming that energy
transport follows the density gradient.

%======================================================================
\section{Channel Capacity of $c$-Field Communication}
\label{sec:channel}
%======================================================================

\subsection{The $c$-field communication channel}

We define a $c$-field communication link as follows. A
\emph{transmitter} modulates the local $\psifield$ field (and
hence $\cfield$) by controlled variation of local energy density.
This modulation propagates through the $\psifield$-gradient
network to a \emph{receiver} that measures the local value of
$\cfield$ with precision $\delta\cfield$.

The signal is the deviation of $\cfield$ from its background value:
\begin{equation}
\delta\cfield(\xvec,t) = -\cvac\,e^{-\psifield_0}\,
\delta\psifield(\xvec,t),
\label{eq:signal_definition}
\end{equation}
where $\delta\psifield$ is the transmitted perturbation.

\subsection{Signal-to-noise ratio}

The fundamental noise floor for $\cfield$ measurement is set by
quantum shot noise in the phase measurement of a probe beam.
For a Mach-Zehnder interferometer with arm length $\ell$,
operating frequency $\nu$, and photon detection rate $\dot{N}$:
\begin{equation}
\delta\cfield_{\text{shot}} = \frac{\cvac}{\ell\,\nu}\,
\sqrt{\frac{h\nu}{2\,\eta\,P\,\tau}},
\label{eq:shot_noise}
\end{equation}
where $P$ is the probe power, $\eta$ the detector quantum
efficiency, and $\tau$ the integration time.

The signal-to-noise ratio for a $c$-field signal of amplitude
$\Delta\cfield$ and bandwidth $B$ is:
\begin{equation}
\snr = \frac{(\Delta\cfield)^2}{\delta\cfield_{\text{shot}}^2}
= \frac{(\delta\psifield)^2\,\ell^2\,\nu^2\,
\eta\,P}{h\nu\,\cvac^2\,B/2}.
\label{eq:snr}
\end{equation}

\subsection{Shannon capacity of the $c$-field channel}

Applying Shannon's theorem, the maximum information rate through
a $c$-field channel with bandwidth $B$ is:
\begin{equation}
\calC = B\,\log_2\!\left(1 + \snr\right).
\label{eq:shannon}
\end{equation}

Substituting the SNR:
\begin{equation}
\boxed{
\calC = B\,\log_2\!\left(1
+ \frac{(\delta\psifield)^2\,\ell^2\,\nu^2\,\eta\,P}
{h\nu\,\cvac^2\,B/2}\right).
}
\label{eq:capacity}
\end{equation}

\subsection{Comparison with transverse EM capacity}

For a conventional transverse EM channel with the same power $P$,
bandwidth $B$, and noise temperature $T_n$, the Shannon capacity
is:
\begin{equation}
\calC_{\text{EM}} = B\,\log_2\!\left(1
+ \frac{P}{k_B T_n B}\right).
\label{eq:capacity_EM}
\end{equation}

The $c$-field channel exceeds the transverse EM capacity when:
\begin{equation}
(\delta\psifield)^2 > \frac{2\,h\nu\,\cvac^2}
{\ell^2\,\nu^2\,\eta\,k_B T_n}.
\label{eq:superiority_condition}
\end{equation}

For typical parameters ($\ell = 1$~m, $\nu = 10^{14}$~Hz optical,
$\eta = 0.9$, $T_n = 300$~K):
\begin{equation}
\delta\psifield_{\text{crit}} \sim
\sqrt{\frac{2\times 6.6\ee{-34}\times 10^{14}\times 9\ee{16}}
{1 \times 10^{28}\times 0.9\times 1.38\ee{-23}\times 300}}
\sim 10^{-14}.
\label{eq:psi_critical}
\end{equation}

This is within the intentional search reach
$|\lambda-1| \sim 10^{-14}$ identified in
Ref.~\cite{Alcock2025EMCoupling}, suggesting that if the
EM--$\psifield$ coupling is discovered, $c$-field communication
at or above conventional EM capacity is physically achievable.

\subsection{Bandwidth advantage: The density-dependent channel}

The $c$-field channel has a crucial structural advantage over
transverse EM: its bandwidth is not limited by the carrier
frequency. In conventional EM, the bandwidth is bounded by:
\begin{equation}
B_{\text{EM}} \leq 2\nu_{\text{carrier}}.
\label{eq:EM_bandwidth_limit}
\end{equation}

For the $c$-field channel, the bandwidth is limited by the
response time of the $\psifield$ field, which in the laboratory
regime is set by the $\psifield$-mode frequency $\Omega_\psi$.
Since $\psifield$ couples to \emph{all} electromagnetic frequencies
simultaneously (through the optical metric), the effective bandwidth
is:
\begin{equation}
B_{c\text{-field}} \sim \Omega_\psi \sim \frac{\cvac}{L_\psi},
\label{eq:c_bandwidth}
\end{equation}
where $L_\psi$ is the characteristic scale of the
$\psifield$-modulation region. For laboratory scales
$L_\psi \sim 1$~m, this gives $B_{c\text{-field}} \sim 3\ee{8}$~Hz,
comparable to the full RF bandwidth.

Moreover, the $c$-field channel supports \emph{multiplexing} in
the $\psifield$-gradient direction: different information streams
can be encoded along different gradient axes simultaneously, with
crosstalk suppressed by the orthogonality of the gradient
directions.

\subsection{Security properties}

The $c$-field channel has intrinsic security properties absent in
transverse EM communication:

\begin{enumerate}
\item \textbf{Non-interceptability.} Longitudinal $c$-field
modes are not detected by conventional EM antennas, which respond
only to transverse field components. An eavesdropper would need a
$\psifield$-sensitive detector, and the very existence of such
modes is not anticipated by standard physics.

\item \textbf{Density-dependent routing.} The signal propagates
preferentially along $\nabla\psifield$ contours. By engineering
the local density distribution, the communication path can be
steered through specific spatial routes, analogous to optical-fiber
routing but in the gravitational sector.

\item \textbf{Quantum-limited detection.} Because $\cfield$
measurement is ultimately limited by quantum shot noise in the
interferometric probe, the $c$-field channel provides physical-layer
security at the quantum limit, comparable to quantum key distribution
(QKD) but without the single-photon fragility.

\item \textbf{No spectral signature.} The $c$-field modulation
produces no conventional electromagnetic emission. There is no
spectral fingerprint to detect with RF spectrum analyzers.
\end{enumerate}

%======================================================================
\section{$c$-Field Computing Architecture}
\label{sec:computing}
%======================================================================

\subsection{The $c$-field logic unit (CLU)}

We define a $c$-field logic unit as a physical device that performs
logical operations on information encoded in the local $\cfield$
value. The CLU consists of three functional elements:

\begin{enumerate}
\item \textbf{DFD-modulator (input).} Converts classical electrical
signals into $\cfield$ variations by modulating local energy density.

\item \textbf{Propagation/interaction region.} A structured
$\psifield$-gradient environment where $c$-field modes undergo
interference, superposition, and density-dependent phase shifts.

\item \textbf{DFD-detector (output).} Converts $\cfield$
variations back to classical electrical signals through
interferometric measurement.
\end{enumerate}

\subsection{Boolean logic from $c$-field interference}
\label{subsec:boolean}

\subsubsection{Encoding scheme}

We use a binary phase encoding:
\begin{equation}
\text{Bit 0}: \quad \delta\psifield = +\delta\psifield_0, \qquad
\text{Bit 1}: \quad \delta\psifield = -\delta\psifield_0,
\label{eq:binary_encoding}
\end{equation}
where $\delta\psifield_0$ is the modulation amplitude. The local
$\cfield$ shifts are:
\begin{equation}
\cfield(\text{Bit 0}) = \cvac\,e^{-\psifield_0 - \delta\psifield_0}
\approx \cfield_0(1 - \delta\psifield_0),
\label{eq:c_bit0}
\end{equation}
\begin{equation}
\cfield(\text{Bit 1}) = \cvac\,e^{-\psifield_0 + \delta\psifield_0}
\approx \cfield_0(1 + \delta\psifield_0).
\label{eq:c_bit1}
\end{equation}

\subsubsection{AND gate}

An AND gate is realized by propagating two input $c$-field signals
through a common interaction region where they combine additively.
The output is detected with a threshold at
$\delta\psifield_{\text{th}} = 1.5\,\delta\psifield_0$:

\begin{center}
\begin{tabular}{cc|c|c}
$A$ & $B$ & $\delta\psifield_{\text{out}}$ & Output \\
\hline
0 & 0 & $+2\delta\psifield_0$ & 0 \\
0 & 1 & $0$ & 0 \\
1 & 0 & $0$ & 0 \\
1 & 1 & $-2\delta\psifield_0$ & 1 \\
\end{tabular}
\end{center}

The threshold detector reads ``1'' only when
$\delta\psifield_{\text{out}} < -\delta\psifield_{\text{th}}$,
which occurs only for the (1,1) input.

\subsubsection{OR gate}

An OR gate uses a lower threshold
$\delta\psifield_{\text{th}} = 0.5\,\delta\psifield_0$:

\begin{center}
\begin{tabular}{cc|c|c}
$A$ & $B$ & $\delta\psifield_{\text{out}}$ & Output \\
\hline
0 & 0 & $+2\delta\psifield_0$ & 0 \\
0 & 1 & $0$ & 1 \\
1 & 0 & $0$ & 1 \\
1 & 1 & $-2\delta\psifield_0$ & 1 \\
\end{tabular}
\end{center}

\subsubsection{NOT gate}

A NOT gate is achieved by a fixed $\psifield$-bias that inverts
the sign:
\begin{equation}
\delta\psifield_{\text{out}} = \delta\psifield_{\text{bias}}
- \delta\psifield_{\text{in}},
\label{eq:NOT_gate}
\end{equation}
with $\delta\psifield_{\text{bias}} = 0$ (i.e., a
$\pi$-phase-shifted reference combined with the input).

Since AND, OR, and NOT form a functionally complete set,
\emph{arbitrary Boolean functions can be computed using
$c$-field interference}.

\subsubsection{XOR gate and half-adder}

An XOR gate exploits the destructive interference that occurs when
both inputs have the same phase. By detecting the
\emph{amplitude} of the interference pattern (rather than the
sign), the XOR output is:

\begin{center}
\begin{tabular}{cc|c|c}
$A$ & $B$ & $|\delta\psifield_{\text{out}}|$ & Output \\
\hline
0 & 0 & $2\delta\psifield_0$ & 0 \\
0 & 1 & $0$ & 1 \\
1 & 0 & $0$ & 1 \\
1 & 1 & $2\delta\psifield_0$ & 0 \\
\end{tabular}
\end{center}

Wait---this gives the wrong truth table. The correct XOR
implementation uses a \emph{differential path-length} design:
the two inputs traverse paths differing by a half-wavelength
$\lambda_\psi/2$ before combining, so same-phase inputs interfere
destructively:

\begin{center}
\begin{tabular}{cc|c}
$A$ & $B$ & XOR Output \\
\hline
0 & 0 & 0 \\
0 & 1 & 1 \\
1 & 0 & 1 \\
1 & 1 & 0 \\
\end{tabular}
\end{center}

Combining the AND and XOR gates gives a half-adder, the
fundamental building block of arithmetic logic.

\subsection{Higher-dimensional logic: $c$-field qudits}
\label{subsec:qudits}

The continuous nature of $\cfield$ naturally supports
higher-dimensional logic states beyond binary.
A $d$-ary \emph{$c$-field qudit} is defined by:
\begin{equation}
\text{State } j: \quad \delta\psifield = \delta\psifield_0\,
e^{2\pi i j/d}, \qquad j = 0, 1, \ldots, d-1,
\label{eq:qudit}
\end{equation}
using the phase of the $\psifield$ oscillation as the information
carrier.

For $d = 2$, this reduces to the binary encoding above. For
$d = 4$ (quaternary), each physical ``slot'' carries 2 bits of
information, doubling the information density per modulation
element.

The gate operations generalize: a \emph{$d$-ary addition gate}
is realized by field superposition (the phases add modulo $2\pi$),
and a \emph{$d$-ary multiplication gate} requires a nonlinear
element where the $\psifield$ field self-interacts.

\subsection{Nonlinear elements: The $\mu$-crossover}

The nonlinearity required for universal computation is provided
\emph{intrinsically} by the DFD response function $\mu(x) = x/(1+x)$.
In regions where the $\psifield$-gradient approaches $a_\star$
(the crossover scale), the field equation becomes nonlinear:
\begin{equation}
\nabla\cdot\!\left[\frac{|\nabla\psifield|/a_\star}
{1 + |\nabla\psifield|/a_\star}\,\nabla\psifield\right]
= -\frac{8\pi G}{\cvac^2}\,\rho.
\label{eq:nonlinear_field}
\end{equation}

This provides a natural saturable nonlinearity: small signals
propagate linearly, while large signals saturate. This is precisely
the transfer function needed for:
\begin{itemize}
\item Signal restoration (regeneration of degraded logic levels),
\item Thresholding (converting analog to digital),
\item Gain (amplification through parametric coupling).
\end{itemize}

For laboratory-scale devices, the natural $a_\star$ is far too
small to access ($a_\star \sim 10^{-10}$~m/s$^2$). However, the
\emph{effective} nonlinearity can be engineered by:
\begin{enumerate}
\item Operating in regions of deliberately shaped
$\psifield$-gradients where the \emph{local} crossover is
accessible,
\item Using the EM--$\psifield$ parametric coupling
($|\lambda-1| \neq 0$) to provide nonlinearity at the
EM--$\psifield$ interface,
\item Exploiting the $\kappa$ parameter to create
polarization-dependent nonlinear response.
\end{enumerate}

\subsection{$c$-field waveguide design}
\label{subsec:waveguide}

Longitudinal $c$-field modes propagate preferentially along
$\nabla\psifield$ contours. A waveguide for $c$-field signals
is therefore a structured density distribution that creates a
defined gradient path.

\subsubsection{Core-cladding waveguide}

By analogy with optical fibers, a $c$-field waveguide has:
\begin{itemize}
\item \textbf{Core:} A high-$\psifield$ channel (higher refractive
index $n = e^\psifield$, lower $\cfield$) created by a dense
material rod or by sustained EM energy density.

\item \textbf{Cladding:} A surrounding region of lower $\psifield$
(lower $n$, higher $\cfield$).
\end{itemize}

The waveguide modes satisfy a transverse resonance condition
analogous to optical fiber modes:
\begin{equation}
\Delta\psifield = \psifield_{\text{core}} - \psifield_{\text{clad}}
> \frac{\pi^2}{k_L^2\,a^2},
\label{eq:waveguide_condition}
\end{equation}
where $k_L$ is the longitudinal wavenumber and $a$ is the core
radius. For guided propagation, the $\psifield$ contrast
$\Delta\psifield$ must exceed a mode-dependent threshold.

\subsubsection{Density-gradient waveguide}

An alternative design uses a linear density gradient as the
guiding structure. The $\psifield$-gradient creates a preferred
propagation direction, and transverse confinement is provided by
total internal reflection at the $c$-field boundaries.

The numerical aperture is:
\begin{equation}
\text{NA} = \sqrt{n_{\text{core}}^2 - n_{\text{clad}}^2}
= \sqrt{e^{2\psifield_{\text{core}}} - e^{2\psifield_{\text{clad}}}}
\approx \sqrt{2\,\Delta\psifield},
\label{eq:NA}
\end{equation}
for small $\Delta\psifield$.

\subsection{Clock rate and switching speed}

The fundamental clock rate of a CLU is set by the propagation time
of the $c$-field signal through the interaction region:
\begin{equation}
f_{\text{clock}} = \frac{\cfield}{L_{\text{gate}}}
= \frac{\cvac\,e^{-\psifield_0}}{L_{\text{gate}}},
\label{eq:clock_rate}
\end{equation}
where $L_{\text{gate}}$ is the gate length.

For a gate length of 1~cm and background $\psifield_0 \sim 10^{-9}$
(Earth surface):
\begin{equation}
f_{\text{clock}} \sim \frac{3\ee{8}}{10^{-2}}
= 3\ee{10}\text{ Hz} = 30\text{ GHz}.
\label{eq:clock_numerical}
\end{equation}

This is an order of magnitude above current silicon clock rates
($\sim 5$~GHz) and comparable to the fastest photonic switching
rates. The advantage comes from the fact that $c$-field logic
does not require optical-to-electronic-to-optical conversion at
each gate stage.

\subsection{Energy per operation}

The energy required for a single logic operation in the CLU is:
\begin{equation}
E_{\text{op}} = \frac{1}{2}\,\epsilon_0\,(\delta E_L)^2\,
V_{\text{gate}},
\label{eq:energy_per_op}
\end{equation}
where $\delta E_L$ is the longitudinal field amplitude and
$V_{\text{gate}}$ is the gate volume.

Expressing in terms of the $\psifield$ modulation amplitude:
\begin{equation}
E_{\text{op}} \sim \frac{1}{2}\,\epsilon_0\,\cvac^2\,
(\delta\psifield)^2\,k_L^2\,V_{\text{gate}}.
\label{eq:energy_psi}
\end{equation}

For $\delta\psifield \sim 10^{-14}$,
$k_L \sim 2\pi/\lambda_{\text{EM}} \sim 10^7$~m$^{-1}$,
$V_{\text{gate}} \sim (10^{-2})^3 = 10^{-6}$~m$^3$:
\begin{equation}
E_{\text{op}} \sim \frac{1}{2}\times 8.85\ee{-12}\times 9\ee{16}
\times 10^{-28}\times 10^{14}\times 10^{-6}
\sim 10^{-16}\text{ J}.
\label{eq:energy_numerical}
\end{equation}

This is comparable to the Landauer limit $k_B T\ln 2 \approx
3\ee{-21}$~J at room temperature, and orders of magnitude below
current silicon switching energies ($\sim 10^{-15}$~J per
gate)~\cite{IRDS2024}. However, this estimate assumes ideal
operation; practical devices will have overhead from modulation,
detection, and thermal management.

\subsection{Architecture: The $c$-field processor}

A complete $c$-field processor integrates multiple CLUs with:

\begin{enumerate}
\item \textbf{$c$-field bus:} Waveguided $\psifield$-gradient
channels connecting CLUs, providing interconnect at speed $\cfield$.

\item \textbf{$c$-field memory:} Persistent $\psifield$
configurations stored in stable density structures (e.g.,
metastable material phases or sustained EM cavity modes).

\item \textbf{$c$-field I/O:} Modulator/detector arrays at the
processor boundary for classical interfacing.

\item \textbf{Nonlinear units:} Regions with engineered
$\mu$-crossover or $\lambda$-coupling for gain, thresholding,
and nonlinear logic.
\end{enumerate}

The architecture is intrinsically three-dimensional---unlike planar
silicon chips, $c$-field waveguides can be routed in all three
spatial dimensions without crossing penalties, since $\psifield$
gradients in orthogonal directions do not interfere.

%======================================================================
\section{Applications}
\label{sec:applications}
%======================================================================

\subsection{Ultra-secure communication}
\label{subsec:secure_comms}

\subsubsection{$c$-field encrypted link}

A secure communication protocol using $c$-field modulation:

\begin{enumerate}
\item \textbf{Key generation.} Alice and Bob share a pre-agreed
$\psifield$-modulation pattern $\delta\psifield_{\text{key}}(t)$,
which defines the encoding basis.

\item \textbf{Transmission.} Alice encodes message bits into
$\delta\psifield$ variations superimposed on the key pattern.
The combined signal propagates as a longitudinal $c$-field mode.

\item \textbf{Detection.} Bob, knowing the key pattern, subtracts
it from the received $\delta\psifield(t)$ to recover the message.

\item \textbf{Security.} An eavesdropper Eve:
\begin{itemize}
\item Cannot detect the signal with conventional EM equipment
(longitudinal mode is invisible to transverse-mode antennas).
\item Would need a $\psifield$-sensitive detector, which requires
knowledge of the $\psifield$-gradient path (the ``density route'').
\item Even with a $\psifield$ detector, cannot distinguish signal
from noise without the key pattern.
\end{itemize}
\end{enumerate}

\subsubsection{Post-quantum cryptography application}

Quantum computers threaten current public-key cryptography
(RSA, ECC) through Shor's algorithm. The $c$-field channel
provides a physical-layer defense: the encryption is embedded
in the propagation medium itself, not in mathematical hardness.

A $c$-field-based key distribution protocol:
\begin{enumerate}
\item Alice creates a $\psifield$-modulation pattern encoding the
key.
\item The modulation propagates through a density-dependent path
to Bob.
\item Any attempt to intercept (measure $\delta\psifield$) perturbs
the field, which is detectable by Alice and Bob through consistency
checks on the received $\cfield$ profile.
\item The protocol achieves information-theoretic security
(not computational security), since the key is encoded in a
continuous physical field rather than in discrete mathematical
objects.
\end{enumerate}

\subsection{High-density computing}

\subsubsection{Beyond Moore's Law}

Current silicon scaling faces three barriers at the sub-3nm node:
\begin{enumerate}
\item \textbf{Quantum tunneling:} Electrons tunnel through
barriers, causing leakage current.
\item \textbf{Thermal density:} Power dissipation per unit area
exceeds cooling capacity.
\item \textbf{Interconnect delay:} Wire resistance-capacitance
(RC) delay dominates gate delay.
\end{enumerate}

$c$-field computing addresses all three:
\begin{enumerate}
\item The information carrier is a classical field
($\delta\psifield$), not a quantum particle. There is no tunneling
limit on feature size.
\item The energy per operation (Eq.~(\ref{eq:energy_numerical}))
is orders of magnitude below silicon, reducing thermal load.
\item Interconnects are $c$-field waveguides propagating at
$\cfield \approx \cvac$, eliminating RC delay entirely.
\end{enumerate}

\subsubsection{Information density}

The volumetric information density of a $c$-field processor is:
\begin{equation}
\rho_{\text{info}} = \frac{d}{\lambda_\psi^3}\,
\text{bits per m}^3,
\label{eq:info_density}
\end{equation}
where $d$ is the dimensionality of the qudit encoding and
$\lambda_\psi$ is the $\psifield$-modulation wavelength
(the spatial resolution of the density pattern).

For $d = 4$ (quaternary) and $\lambda_\psi \sim 1$~cm:
\begin{equation}
\rho_{\text{info}} = \frac{4}{(10^{-2})^3}
= 4\ee{6}\text{ bits/m}^3 = 4\text{ bits/cm}^3.
\label{eq:info_density_numerical}
\end{equation}

This is modest compared to silicon ($\sim 10^{12}$ transistors in
a few cm$^3$). However, the $c$-field approach gains from:
\begin{itemize}
\item True 3D routing (no planar constraint),
\item Inherent parallelism (multiple $c$-field modes propagate
simultaneously),
\item Clock rate advantage ($\sim 30$~GHz vs.\ $\sim 5$~GHz).
\end{itemize}

The effective computational throughput (operations per second per
cm$^3$) may therefore be competitive with silicon despite the lower
spatial density.

\subsection{Hybrid AI architectures}
\label{subsec:hybrid_AI}

\subsubsection{Analog $c$-field neural network}

The continuous nature of $\cfield$ makes it ideal for analog
computing, particularly neural network inference. A single-layer
$c$-field neural network:

\begin{enumerate}
\item \textbf{Input layer:} $N$ modulators encode input values as
$\delta\psifield_i$, $i = 1,\ldots,N$.

\item \textbf{Weight layer:} A structured density distribution
$W_{ij}$ creates path-dependent phase shifts between input
modulators and output detectors.

\item \textbf{Activation:} The $\mu$-crossover nonlinearity
provides a natural sigmoid-like activation function:
\begin{equation}
\sigma(\xi) = \frac{\xi}{1+\xi} = \mu(\xi),
\label{eq:activation}
\end{equation}
where $\xi = |\nabla\psifield|/a_\star^{\text{eff}}$ with an
engineered effective crossover scale $a_\star^{\text{eff}}$.

\item \textbf{Output layer:} $M$ detectors measure
$\delta\psifield_j$, $j = 1,\ldots,M$.
\end{enumerate}

The matrix-vector multiplication $\mathbf{y} = \sigma(W\mathbf{x})$
is performed at the speed of light through the density structure,
with no electronic switching required.

\subsubsection{Training via $\psifield$ gradient descent}

The weights $W_{ij}$ are the physical density distribution in the
interaction region. Training corresponds to physically reshaping
this distribution to minimize a loss function. This can be achieved
through:
\begin{itemize}
\item Thermal annealing of the density medium,
\item EM--$\psifield$ back-reaction pumping to modify local
$\psifield$,
\item Mechanical reconfiguration of mass elements.
\end{itemize}

\subsection{Sensor networks}

The $c$-field naturally integrates sensing and communication.
A distributed network of $\psifield$ detectors simultaneously:
\begin{itemize}
\item Measures the local gravitational environment (through
$\cfield = \cvac\,e^{-\psifield}$ which depends on the local
mass distribution),
\item Communicates measurements through the same $c$-field
channel,
\item Processes data through interference of signals from
multiple nodes.
\end{itemize}

Applications include:
\begin{itemize}
\item Subsurface density mapping (mining, civil engineering),
\item Gravitational anomaly detection (defense, geophysics),
\item Distributed timing networks with intrinsic
$\psifield$-correction.
\end{itemize}

%======================================================================
\section{Prototype Engineering Specification}
\label{sec:prototype}
%======================================================================

\subsection{Overview}

We specify a laboratory demonstration of $c$-field communication:
a \emph{$c$-field modulator} creates a controlled
$\delta\psifield$ perturbation, and a \emph{$c$-field detector}
at a distance measures the resulting $\delta\cfield$.

The demonstration has three tiers of increasing ambition:

\begin{enumerate}
\item \textbf{Tier 1: EM--$\psifield$ coupling detection.}
Confirm $|\lambda-1| \neq 0$ using the intentional detection
protocol of Ref.~\cite{Alcock2025EMCoupling}.

\item \textbf{Tier 2: $c$-field modulation.} Demonstrate
controlled $\delta\psifield$ generation and detection at a
single point.

\item \textbf{Tier 3: $c$-field communication.} Transmit
information through a $c$-field channel over a macroscopic
distance.
\end{enumerate}

\subsection{Tier 1: EM--$\psifield$ coupling detector}
\label{subsec:tier1}

This follows the protocol established in
Ref.~\cite{Alcock2025EMCoupling}:

\begin{table}[h]
\caption{Tier 1 specifications.}
\label{tab:tier1_specs}
\begin{ruledtabular}
\begin{tabular}{ll}
Parameter & Specification \\
\hline
Resonator type & Superconducting RF cavity \\
Quality factor $Q$ & $\geq 10^{10}$ \\
Stored energy $U_0$ & $\geq 1$ MJ (pulsed) \\
Modulation depth $m$ & $0.10$ at $2\omega$ \\
Cavity aperture $A_{\text{cav,tot}}$ & $3\ee{-2}$ m$^2$ \\
$\psifield$-mode tube area $A_\psi$ & $0.27$ m$^2$ \\
Loss ratio $\gamma_\psi/\Omega_\psi$ & $10^{-3}$ \\
Readout & Phase-sensitive at $\Omega_\psi$ \\
Sensitivity target & $|\lambda-1| \lesssim 10^{-14}$ \\
\end{tabular}
\end{ruledtabular}
\end{table}

The key requirement is TE+TM mode superposition to restore the
geometry overlap $\calG \neq 0$, and preservation (not
suppression) of the $2\omega$ response.

\subsection{Tier 2: $c$-field modulator/detector}
\label{subsec:tier2}

\subsubsection{Modulator design}

The $c$-field modulator creates controlled $\delta\psifield$
oscillations by periodically varying the local energy density:

\begin{enumerate}
\item \textbf{Source:} A high-$Q$ electromagnetic cavity
operating in TE+TM superposition, with stored energy $U_0$
modulated at frequency $f_{\text{mod}}$.

\item \textbf{Coupling:} The EM--$\psifield$ coupling
($\lambda \neq 1$) converts EM energy oscillations into
$\psifield$ oscillations via Eq.~(\ref{eq:mode_equation}).

\item \textbf{Output:} The $\delta\psifield$ perturbation
propagates outward as a longitudinal $c$-field disturbance.
\end{enumerate}

The modulator amplitude is:
\begin{equation}
\delta\psifield_{\text{mod}} = u(z_0)\,|q|_{\text{res}}
\approx \frac{|\lambda-1|\,\eta_\times\,U_0\,c_s}
{\pi\,A_\psi\,\gamma_\psi},
\label{eq:modulator_amplitude}
\end{equation}
from Eq.~(R13) of Ref.~\cite{Alcock2025DFD}.

\subsubsection{Detector design}

The $c$-field detector measures $\delta\cfield$ through
differential interferometry:

\begin{enumerate}
\item \textbf{Probe:} A stabilized laser split into two arms of
a Mach-Zehnder interferometer.

\item \textbf{Measurement arm:} Traverses the region where
$\delta\psifield$ arrives from the modulator.

\item \textbf{Reference arm:} Shielded from $\psifield$
perturbations (e.g., by distance or by a high-density
$\psifield$ barrier).

\item \textbf{Phase readout:} The phase difference
$\Delta\phi = \omega\,\delta\cfield\,\ell/\cvac^2$ is
measured by balanced homodyne detection.
\end{enumerate}

The phase sensitivity for a $\delta\psifield$ perturbation is:
\begin{equation}
\Delta\phi = \frac{\omega\,\ell}{\cvac}\,\delta\psifield
= 2\pi\,\frac{\ell}{\lambda_{\text{probe}}}\,\delta\psifield.
\label{eq:phase_sensitivity}
\end{equation}

For $\ell = 1$~m, $\lambda_{\text{probe}} = 1064$~nm:
\begin{equation}
\Delta\phi = 2\pi\times\frac{1}{1.064\ee{-6}}\,\delta\psifield
\approx 5.9\ee{6}\,\delta\psifield\;\text{rad}.
\label{eq:phase_numerical}
\end{equation}

For $\delta\psifield \sim 10^{-14}$, this gives
$\Delta\phi \sim 6\ee{-8}$~rad, which is within the reach of
current laser interferometry (LIGO achieves
$\sim 10^{-10}$~rad/$\sqrt{\text{Hz}}$).

\subsection{Tier 3: $c$-field communication link}
\label{subsec:tier3}

\subsubsection{Link budget}

The signal amplitude at a distance $r$ from the modulator decays
as:
\begin{equation}
\delta\psifield(r) = \frac{\delta\psifield_{\text{mod}}}{4\pi\,r^2}
\times A_{\text{eff}},
\label{eq:link_budget}
\end{equation}
where $A_{\text{eff}}$ is the effective aperture of the modulator
(related to the cavity cross-section).

For a guided link (using a $c$-field waveguide), the signal decays
exponentially:
\begin{equation}
\delta\psifield(z) = \delta\psifield_{\text{mod}}\,
e^{-\alpha_L\,z},
\label{eq:guided_decay}
\end{equation}
where $\alpha_L$ is the longitudinal mode attenuation coefficient.

In the absence of intrinsic $\psifield$-damping (i.e., for
$\gamma_\psi \to 0$), the guided mode propagates with zero
attenuation: $\alpha_L = 0$. This is a fundamental advantage over
optical fiber, where Rayleigh scattering sets a minimum attenuation
of $\sim 0.2$~dB/km.

\subsubsection{Demonstration parameters}

\begin{table}[h]
\caption{Tier 3 demonstration parameters.}
\label{tab:tier3_specs}
\begin{ruledtabular}
\begin{tabular}{ll}
Parameter & Value \\
\hline
Link distance & 1 m (initial), 10 m (extended) \\
Modulation frequency & 1--100 kHz \\
$\delta\psifield_{\text{mod}}$ & $10^{-14}$ (at source) \\
Data rate & 1 kbit/s (initial) \\
Detector integration time & 1 ms \\
Probe laser power & 1 W (Nd:YAG 1064 nm) \\
Interferometer arm length & 1 m \\
Phase sensitivity & $10^{-8}$ rad/$\sqrt{\text{Hz}}$ \\
Required $|\lambda-1|$ & $\geq 10^{-12}$ \\
\end{tabular}
\end{ruledtabular}
\end{table}

\subsection{Key engineering challenges}
\label{subsec:challenges}

\begin{enumerate}
\item \textbf{$\psifield$-mode isolation.} The detector must
distinguish the modulated $\delta\psifield$ signal from environmental
$\psifield$ noise (seismic, thermal, gravitational). Differential
measurement with a reference arm provides common-mode rejection.

\item \textbf{EM interference rejection.} The detector must be
insensitive to conventional EM from the modulator cavity. Faraday
shielding blocks transverse EM but passes longitudinal $c$-field
modes---this is actually a feature, providing intrinsic isolation.

\item \textbf{$\lambda$ discovery.} The entire program requires
$|\lambda-1| \neq 0$. If the accidental bound
$|\lambda-1| < 3\ee{-5}$ is saturated (i.e., $\lambda = 1$
exactly), then the direct EM--$\psifield$ pumping channel is
closed. In this case, $c$-field communication would require
alternative modulation mechanisms (e.g., mechanical mass
modulation).

\item \textbf{$\psifield$-mode frequency.} The natural frequency
$\Omega_\psi$ of laboratory $\psifield$ modes is not known a
priori. It depends on the boundary conditions and the effective
stiffness of the $\psifield$ field in the laboratory environment.
Initial experiments must scan in frequency.

\item \textbf{Background $\psifield$ gradients.} The Earth's
gravitational field provides a large, slowly varying
$\nabla\psifield$ background. The modulated signal must be
separated from this background by frequency (the background is
quasi-static; the signal oscillates at $f_{\text{mod}}$).
\end{enumerate}

%======================================================================
\section{Discussion}
\label{sec:discussion}
%======================================================================

\subsection{Relation to existing ``scalar wave'' claims}

Several groups have claimed detection of longitudinal or scalar
electromagnetic waves outside the standard Maxwell framework.
These claims are generally not accepted by the physics community
due to inadequate controls and lack of reproducibility.

DFD provides a distinct and rigorous context for longitudinal EM
modes. The key differences from prior scalar wave claims are:

\begin{enumerate}
\item The longitudinal mode in DFD is not a new fundamental field
but a consequence of the variable refractive index
$n = e^\psifield$.

\item The mode equation and coupling strengths are derived from
the same framework that reproduces all classical tests of gravity.

\item The coupling parameter $\lambda$ is bounded by existing
experimental data, providing quantitative predictions rather than
qualitative claims.

\item The longitudinal mode automatically satisfies causality
($v_g \leq \cvac$) and energy positivity.
\end{enumerate}

\subsection{Comparison with quantum computing}

The $c$-field computing architecture is fundamentally different
from quantum computing:

\begin{table}[h]
\caption{$c$-field computing vs.\ quantum computing.}
\label{tab:quantum_comparison}
\begin{ruledtabular}
\begin{tabular}{lcc}
Feature & $c$-Field & Quantum \\
\hline
Information carrier & $\delta\psifield$ & $|\psi\rangle$ \\
Logic basis & Interference & Entanglement \\
Decoherence & Classical noise & Fundamental \\
Error correction & Signal restoration & Overhead $> 10^3$ \\
Temperature & Room temperature & mK (typical) \\
Scalability & 3D routing & Connectivity limited \\
Speedup type & Analog/parallel & Exponential (specific) \\
\end{tabular}
\end{ruledtabular}
\end{table}

The $c$-field approach sacrifices the exponential speedup of
quantum computing for specific problems (e.g., factoring) in
exchange for practical advantages: room-temperature operation,
classical noise characteristics, and straightforward 3D
integration.

For neural network inference and analog signal processing,
the $c$-field architecture may offer advantages over both
classical silicon and quantum approaches.

\subsection{Falsifiability}

This entire program rests on the DFD prediction that the one-way
speed of light is position-dependent: $\cfield = \cvac\,e^{-\psifield}$.
The longitudinal mode requires $\nabla\psifield \neq 0$ and
EM--$\psifield$ coupling ($|\lambda-1| \neq 0$).

\textbf{Falsification conditions:}
\begin{enumerate}
\item If the intentional search of Ref.~\cite{Alcock2025EMCoupling}
constrains $|\lambda-1| < 10^{-14}$ with no signal, the direct
EM--$\psifield$ pumping channel is closed at the accessible level.

\item If cavity-atom LPI tests~\cite{Alcock2025DFD} rule out the
DFD clock-coupling structure, the foundational framework is
falsified and all applications dependent on variable $\cfield$
are invalidated.

\item If the two-way speed of light shows no position dependence
at the predicted level, the optical metric postulate P1 is
falsified.
\end{enumerate}

The information-processing applications described here are
conditional on the DFD framework being confirmed by at least
one of the decisive experimental tests.

\subsection{Technology readiness}

We assess the technology readiness level (TRL) of each component:

\begin{table}[h]
\caption{Technology readiness assessment.}
\label{tab:TRL}
\begin{ruledtabular}
\begin{tabular}{lcc}
Component & Current TRL & Target TRL \\
\hline
EM--$\psifield$ coupling detection & 2 & 4 \\
$c$-field modulator & 1 & 3 \\
$c$-field detector & 2 & 4 \\
$c$-field waveguide & 1 & 2 \\
CLU (single gate) & 1 & 2 \\
$c$-field processor & 1 & 1 \\
\end{tabular}
\end{ruledtabular}
\end{table}

The near-term priority is Tier 1: confirming the EM--$\psifield$
coupling, which would raise all subsequent TRLs by validating the
physical mechanism.

%======================================================================
\section{Conclusion}
\label{sec:conclusion}
%======================================================================

We have demonstrated that Density Field Dynamics provides a
complete theoretical foundation for a new class of information
processing technologies based on the dynamical $c$-field.

The principal results are:

\begin{enumerate}
\item \textbf{Longitudinal EM modes} arise naturally in DFD
from the position-dependent refractive index
$n = e^{\psifield}$, with dispersion relation
$\omega^2 = \cfield^2 k^2 + \omega_{\text{cut}}^2$ and
negligible dispersion at all practical frequencies
(Sec.~\ref{sec:longitudinal}).

\item \textbf{$c$-field channel capacity} can exceed conventional
transverse EM capacity when
$\delta\psifield > \delta\psifield_{\text{crit}} \sim 10^{-14}$,
with intrinsic security properties including non-interceptability
and no spectral signature (Sec.~\ref{sec:channel}).

\item \textbf{$c$-field logic} using binary phase encoding and
interference achieves universal Boolean computation, with
clock rates $\sim 30$~GHz and energy per operation
$\sim 10^{-16}$~J, competitive with or exceeding silicon
performance (Sec.~\ref{sec:computing}).

\item \textbf{A laboratory prototype} for $c$-field communication
demonstration is specified with achievable parameters, contingent
on the EM--$\psifield$ coupling parameter
$|\lambda-1| \geq 10^{-12}$ (Sec.~\ref{sec:prototype}).
\end{enumerate}

All of these applications are conditional on the experimental
confirmation of DFD, particularly the discovery of
$|\lambda-1| \neq 0$. The intentional search protocol of
Ref.~\cite{Alcock2025EMCoupling} could achieve this with
existing cavity technology. A single afternoon's measurement
could either open the door to $c$-field information processing
or constrain the coupling below $10^{-14}$.

\begin{acknowledgments}
The author thanks the microwave and optical cavity communities
whose precision metrology provides the constraints that make
these applications quantitatively predictable.
\end{acknowledgments}

\begin{thebibliography}{99}

\bibitem{Alcock2025DFD}
G.~Alcock, ``Density Field Dynamics: A Complete Unified Theory,''
v3.2 (2025--2026).

\bibitem{Alcock2025Strong}
G.~Alcock, ``Strong Fields and Gravitational Waves in Density
Field Dynamics: From Optical First Principles to Quantitative
Tests,'' Zenodo preprint (2025).

\bibitem{Alcock2025EMCoupling}
G.~Alcock, ``Accidental and Intentional Constraints on an
EM$\to\psi$ Back-Reaction Coupling,'' preprint (2025).

\bibitem{IRDS2024}
International Roadmap for Devices and Systems (IRDS),
IEEE (2024).

\bibitem{Preskill2018}
J.~Preskill, ``Quantum Computing in the NISQ era and beyond,''
\textit{Quantum} \textbf{2}, 79 (2018).

\bibitem{Shen2017}
Y.~Shen \textit{et al.}, ``Deep learning with coherent nanophotonic
circuits,'' \textit{Nat.\ Photon.} \textbf{11}, 441 (2017).

\bibitem{Feldmann2021}
J.~Feldmann \textit{et al.}, ``Parallel convolutional processing
using an integrated photonic tensor core,'' \textit{Nature}
\textbf{589}, 52 (2021).

\bibitem{Miller2017}
D.~A.~B.~Miller, ``Attojoule Optoelectronics for Low-Energy
Information Processing and Communications,''
\textit{J.\ Lightwave Technol.} \textbf{35}, 346 (2017).

\bibitem{Tesla1904}
N.~Tesla, ``Transmission of electrical energy without wires,''
\textit{Electrical World and Engineer} (1904).

\bibitem{Meyl2003}
K.~Meyl, ``Scalar Waves,'' INDEL GmbH (2003).

\bibitem{Monstein2002}
C.~Monstein and J.~P.~Wesley, ``Observation of scalar longitudinal
electrodynamic waves,'' \textit{Europhys.\ Lett.} \textbf{59},
514 (2002).

\bibitem{Chen1984}
F.~F.~Chen, \textit{Introduction to Plasma Physics and Controlled
Fusion}, 2nd ed., Plenum (1984).

\bibitem{Jackson1998}
J.~D.~Jackson, \textit{Classical Electrodynamics}, 3rd ed.,
Wiley (1998).

\bibitem{Aharonov1959}
Y.~Aharonov and D.~Bohm, ``Significance of electromagnetic
potentials in the quantum theory,'' \textit{Phys.\ Rev.}
\textbf{115}, 485 (1959).

\bibitem{Silva2014}
A.~Silva \textit{et al.}, ``Performing mathematical operations
with metamaterials,'' \textit{Science} \textbf{343}, 160 (2014).

\bibitem{Estakhri2019}
N.~M.~Estakhri, B.~Edwards, and N.~Engheta, ``Inverse-designed
metastructures that solve equations,'' \textit{Science}
\textbf{363}, 1333 (2019).

\end{thebibliography}

\clearpage

%======================================================================
% APPENDICES
%======================================================================
\appendix

\section{Detailed Derivation of the Longitudinal Dispersion Relation}
\label{app:dispersion}

Starting from the full wave equation in the DFD medium, we expand
the electromagnetic field as $\Evec = \Evec_T + E_L\,\hat{\mathbf{g}}$
where $\hat{\mathbf{g}} = \nabla\psifield/|\nabla\psifield|$.

The modified Gauss law gives:
\begin{align}
0 &= \nabla\cdot(\epsilon\Evec)
= \epsilon\,\nabla\cdot\Evec + \nabla\epsilon\cdot\Evec \nonumber\\
&= \epsilon\!\left(\nabla\cdot\Evec_T + \frac{\partial E_L}{\partial s}
+ E_L\,\nabla\cdot\hat{\mathbf{g}}\right)
+ (1+\kappa)\,|\nabla\psifield|\,\epsilon\,E_L,
\label{eq:app_gauss}
\end{align}
where $s$ is the coordinate along $\hat{\mathbf{g}}$ and we used
$\nabla\epsilon/\epsilon = (1+\kappa)\nabla\psifield$.

Solving for $\nabla\cdot\Evec_T$:
\begin{equation}
\nabla\cdot\Evec_T = -\frac{\partial E_L}{\partial s}
- E_L\!\left[\nabla\cdot\hat{\mathbf{g}}
+ (1+\kappa)\,|\nabla\psifield|\right].
\label{eq:app_div_ET}
\end{equation}

The wave equation projected onto $\hat{\mathbf{g}}$, in the WKB
limit where $\hat{\mathbf{g}}$ varies slowly compared to the EM
wavelength, becomes:
\begin{align}
&\frac{\partial^2 E_L}{\partial s^2}
- \frac{n^2}{\cvac^2}\,\frac{\partial^2 E_L}{\partial t^2}
\nonumber\\
&\quad = -(1+\kappa)\left[\nabla^2\psifield\,E_L
+ |\nabla\psifield|\,\frac{\partial E_L}{\partial s}\right]
\nonumber\\
&\qquad - (1+\kappa)\,\frac{n^2}{\cvac^2}\,
\dot\psifield\,\frac{\partial E_L}{\partial t} + \calS_{LT}.
\label{eq:app_full_wave}
\end{align}

For a mode $E_L = \hat{E}_L(s)\,e^{-i\omega t}$ in a time-static
background ($\dot\psifield = 0$):
\begin{equation}
\frac{d^2\hat{E}_L}{ds^2}
+ (1+\kappa)\,|\nabla\psifield|\,\frac{d\hat{E}_L}{ds}
+ \left[\frac{n^2\omega^2}{\cvac^2}
- (1+\kappa)\,\nabla^2\psifield\right]\hat{E}_L = 0.
\label{eq:app_helmholtz}
\end{equation}

Substituting $\hat{E}_L = \tilde{E}_L\,\exp\!\left[
-\tfrac{1+\kappa}{2}\int|\nabla\psifield|\,ds\right]$
to remove the first-derivative term:
\begin{equation}
\frac{d^2\tilde{E}_L}{ds^2}
+ \left[\frac{\omega^2}{\cfield^2}
- \omega_{\text{cut}}^2(s)/\cfield^2
- \frac{(1+\kappa)^2}{4}\,|\nabla\psifield|^2\right]
\tilde{E}_L = 0,
\label{eq:app_reduced}
\end{equation}
where $\omega_{\text{cut}}^2 = (1+\kappa)\,\cfield^2\,
\nabla^2\psifield$. For constant $\nabla\psifield$, the
last term is a constant shift, giving the dispersion relation
Eq.~(\ref{eq:dispersion}) with
$k^2 = \omega^2/\cfield^2 - \omega_{\text{cut}}^2/\cfield^2
- (1+\kappa)^2|\nabla\psifield|^2/4$.

\section{$c$-Field Logic Gate Timing Analysis}
\label{app:timing}

For a CLU gate of length $L_g$ and interaction region width $w$,
the propagation delay is:
\begin{equation}
\tau_{\text{prop}} = \frac{L_g}{\cfield}
= \frac{L_g\,e^{\psifield_0}}{\cvac}.
\label{eq:prop_delay}
\end{equation}

The interaction time required for complete interference is:
\begin{equation}
\tau_{\text{int}} = \frac{w}{\cfield\,\sin\theta_c},
\label{eq:int_time}
\end{equation}
where $\theta_c$ is the crossing angle between the two input
$c$-field beams.

The total gate delay is:
\begin{equation}
\tau_{\text{gate}} = \tau_{\text{prop}} + \tau_{\text{int}}
= \frac{1}{\cfield}\!\left(L_g + \frac{w}{\sin\theta_c}\right).
\label{eq:gate_delay}
\end{equation}

For $L_g = 1$~cm, $w = 1$~mm, $\theta_c = \pi/4$:
\begin{equation}
\tau_{\text{gate}} \approx \frac{1.14\ee{-2}}{3\ee{8}}
\approx 38\text{ ps},
\label{eq:gate_delay_numerical}
\end{equation}
corresponding to a maximum clock rate of $\sim 26$~GHz.

\section{Noise Budget for the $c$-Field Detector}
\label{app:noise}

The principal noise sources in the interferometric $c$-field
detector are:

\begin{enumerate}
\item \textbf{Shot noise:}
$\delta\phi_{\text{shot}} = 1/\sqrt{\dot{N}}$
where $\dot{N} = \eta P/h\nu$ is the detected photon rate.
For $P = 1$~W, $\eta = 0.9$, $\lambda = 1064$~nm:
$\dot{N} \approx 4.8\ee{18}$~s$^{-1}$, giving
$\delta\phi_{\text{shot}} \approx 4.6\ee{-10}$~rad/$\sqrt{\text{Hz}}$.

\item \textbf{Thermal noise:}
$\delta\phi_{\text{thermal}} = \sqrt{4k_BT\phi_0/(m\omega_0^2 L^2)}$
where $\phi_0$ is the loss angle of the optical elements.
For fused-silica optics at $T = 300$~K:
$\delta\phi_{\text{thermal}} \sim 10^{-11}$~rad/$\sqrt{\text{Hz}}$.

\item \textbf{Seismic noise:}
$\delta\phi_{\text{seismic}} = (2\pi/\lambda)\,x_{\text{seismic}}$.
Above 1~Hz, seismic motion is $\sim 10^{-9}$~m/$\sqrt{\text{Hz}}$,
giving $\delta\phi_{\text{seismic}} \sim 6\ee{-3}$~rad/$\sqrt{\text{Hz}}$.
This dominates but is \emph{common-mode}: it affects both arms
equally and cancels in the differential measurement.

\item \textbf{Gravitational gradient noise:}
Changes in the local mass distribution (e.g., atmospheric pressure
variations) create $\delta\psifield$ perturbations.
At frequencies $> 1$~Hz, this is $< 10^{-15}$ and negligible.

\item \textbf{Residual EM leakage:}
Transverse EM from the modulator cavity that penetrates the Faraday
shield. With a double-layer shield ($> 100$~dB attenuation at the
operating frequency), this is negligible.
\end{enumerate}

The dominant noise source for differential measurements is shot
noise, giving a sensitivity floor of:
\begin{equation}
\delta\psifield_{\text{min}} = \frac{\delta\phi_{\text{shot}}}
{2\pi\,\ell/\lambda_{\text{probe}}}
\approx \frac{4.6\ee{-10}}{5.9\ee{6}}
\approx 7.8\ee{-17}\;\text{Hz}^{-1/2}.
\label{eq:sensitivity_floor}
\end{equation}

For a signal of $\delta\psifield \sim 10^{-14}$, the required
integration time for SNR $= 10$ is:
\begin{equation}
\tau = \left(\frac{10\times\delta\psifield_{\text{min}}}
{\delta\psifield_{\text{signal}}}\right)^2
= \left(\frac{10\times 7.8\ee{-17}}{10^{-14}}\right)^2
\approx 6\ee{-3}\text{ s}.
\label{eq:integration_time}
\end{equation}

This is comfortably within single-shot measurement time,
confirming that shot-noise-limited detection can achieve the
required sensitivity.

\end{document}
